\documentclass[11pt,a4paper]{article}
\usepackage[utf8]{inputenc}
\usepackage[french]{babel}
\usepackage[T1]{fontenc}
\usepackage{amsmath}
\usepackage{amsfonts}
\usepackage{amssymb}
\usepackage{graphicx}
\usepackage{enumitem}

\usepackage{fancybox}
\cornersize{3} 

\usepackage{tikz}
\usetikzlibrary {arrows.meta} 


\usepackage{minted}
\newcommand{\mip}[1]{\mintinline{python}{#1}}
\usepackage{hyperref}
\usepackage{pstricks,pst-plot,pst-3dplot,pst-grad,pst-tree,pst-math,pst-eucl}
\usepackage{pst-plot,pst-tree,pst-eucl,pst-math,pstricks-add}
\usepackage[margin=1cm,bottom=1.5cm]{geometry}
\begin{document}
\fbox{\textsc{25NSIJ2AN1}}


\bigskip

\textbf{Exercice 1} (6 points)

\medskip

\textit{Cet exercice porte sur les tableaux, les dictionnaires, les arbres binaires, la
programmation en Python}

\textit{et la récursivité.}

\smallskip

\textbf{Partie A}

\begin{enumerate}
\item L'écriture binaire du code ASCII 97 du caractère \mip{'a'} est 0110.0001

\item L'appel \mip{replique([0,0,0,1,0,1])} a pour résultat 
\mip{[0,0,0,0,0,0,0,0,0,1,1,1,0,0,0,1,1,1]}.

\item On complète les lignes demandées.

\begin{minted}[linenos,firstnumber=13]{python}
    if x in nb_occ:
        nb_occ[x]+=1
    else;
        nb_occ[x]=1
\end{minted} 


\item On complète la fonction \mip{majorite}.%, en prenant bien entendu bien soin de respecter l'indentation qui est un élément fondamental de la syntaxe Python.

\begin{minted}[linenos,firstnumber=10]{python}
    for cle in dict.keys():
        if dict[cle]>valeur_max:
             cle_max=cle
             valeur_max=dict[cle]
    return cle_max
\end{minted} 

\end{enumerate}


\textbf{Partie B}

\begin{enumerate}[resume]
\item On a entouré le bit qui a subi une erreur.

\vspace{-1.5cm}
\hspace{10cm}\begin{tabular}{|c|c|c|}
\hline
\rule[-2pt]{0pt}{13pt}1&\Ovalbox{1}&1\\
\hline
1&1&0\\
\hline
0&1&1\\
\hline
\end{tabular}


\item On écrit la fonction \mip{erreur_colonne} demandée.

\begin{minted}{python}
def erreur_colonne(mat):
    if (mat[0][0]+mat[1][0]+mat[2][0])%2=1:
        return 0
    elif (mat[0][1]+mat[1][1]+mat[2][1])%2=1:
        return 1
    return 2
\end{minted}

\end{enumerate}



\textbf{Partie C}

\begin{enumerate}[resume]
\item Avec un code reçu 1010000, le code de Hamming le plus proche est 1110000 qui correspond au mot 1000.

\item On complète la fonction \mip{corriger_erreur}

\begin{minted}[linenos,firstnumber=12]{python}
def corriger_erreur(code_recu):
    if code_recu in hamming_4_7:
        return code_recu
    else:
        # Copie du code reçu avec une compréhension
        code=[bit for bit in code_recu]
        for indice in range(7):
            # Inversion du bit d'indice courant
            code[indice]=(code[indice]+1)%2
            if code in hamming_4_7:
                return code
            else:
                # Réinit. du bit d'indice courant
                code[indice]=(code[indice]+1)%2
\end{minted}

\item L'arbre décodeur complet du code de Hamming (4,7) comporte $2^7=128$ feuilles. 

\newpage

\item On complète la fonction récursive \mip{decode}.

\begin{minted}[linenos,firstnumber=12]{python}
    if code[i]==0:
        return decode(arbre.gauche, code, i+1)
    if code[i]==1:
        return decode(arbre.droit, code, i+1)
\end{minted}

\end{enumerate}




\bigskip

\textbf{Exercice 2} (6 points)

\medskip

\textit{Cet exercice porte sur la gestion des bugs, l’algorithmique, les structures de données et la programmation}

\textit{orientée objet.}

\smallskip

\textbf{Partie A}
\begin{enumerate}
\item Pour un collier de 8 bonbons, ceux-ci sont mangés dans l'ordre 0, 3, 6, 2, 7, 5, 1 et le bonbon restant est le 4.

\item L'erreur \mip{NameError} est une erreur indiquant que Python ne connaît pas l'identifiant \mip{true}. La valeur booléenne vraie est notée \mip{True} en Python, Bob doit faire trop de JavaScript, où les valeurs booléennes sont en minuscules.

\item On complète le code de la fonction \mip{dernier}.

\begin{minted}[linenos]{python}
def dernier(n):
    collier=[True]*n
    indice=0
    collier[indice]=False
    for etape in range(n-1):
        nb_bonbons_vus=0
        while nb_bonbons_vus<3:
            indice+=1
            if indice==n:
                indice=0
            if collier[indice]:
                nb_bonbons_vus+=1
        collier[indice]=False
    return indice    
\end{minted}
\end{enumerate}

\vspace{-6.5cm}

\hfill\begin{minipage}{9cm}
\textbf{Partie B}

\begin{enumerate}[resume]
\item Une \mip{File} est une structure LIFO.

\item L'affichage sera le suivant.
\begin{minted}{python}
>>> f.affiche()
(Tête) 3 4 1 2 (Queue)
\end{minted}

\item On écrit une fonction \mip{dernier_file} qui utilise une \mip{File}.
\begin{minted}{python}
def dernier_file(n):
    f=File()
    for bonbon in range(n):
        f.enfile(bonbon)
    while True:
        bonbon=f.defile()
        if f.est_vide():
            return bonbon
        f.enfile(f.defile())
        f.enfile(f.defile())
\end{minted}

\end{enumerate}
\end{minipage}


\vspace{-2cm}


\textbf{Partie C}

\begin{enumerate}
\setcounter{enumi}{6}
\item \mip{pred}, \mip{valeur} et \mip{succ} sont des attributs 

de la classe \mip{Bonbon}.
\item Après exécution du code proposée, \mip{a} a comme valeur \mip{1}

et \mip{b} comme valeur \mip{2}.

\item On complète le code de la fonction \mip{creer_collier}.

\begin{minted}[linenos]{python}
def creer_collier(n):
    premier=Bonbon(0)
    actuel=premier
    for i in range(1,n):
        nouveau=Bonbon(i)
        actuel.succ=nouveau
        nouveau.pred=actuel
        actuel=nouveau
    actuel.succ=premier
    premier.pred=actuel
    return premier
\end{minted}

\item On dessine le collier obtenu. 

Le bonbon \mip{0} a été mangé et n'appartient plus au collier.

\vspace{-3.5cm}
\hfill\begin{tikzpicture}
\foreach \x in {0,1,2,3} {
\draw({3*\x},0)--({3*\x+2},0)--({3*\x+2},1)--({3*\x},1)--cycle;
\draw({3*\x+0.5},0)--({3*\x+0.5},1);
\draw({3*\x+1.5},0)--({3*\x+1.5},1);
\draw({3*\x+1},0.5)node{\x};
}
\draw(4.75,0.7)node{$\bullet$};\draw[-{Latex[length=3mm]},dashed](4.75,0.7)--(6,.7);
\draw(7.75,0.7)node{$\bullet$};\draw[-{Latex[length=3mm]},dashed](7.75,0.7)--(9,.7);
\draw(10.75,0.7)node{$\bullet$};\draw[-{Latex[length=3mm]},dashed](10.75,0.7)--(11.5,0.7)--(11.5,1.5)--(4,1.5)--(4,1);

\draw(6.25,0.3)node{$\bullet$};\draw[-{Latex[length=3mm]},dashed](6.25,0.3)--(5,.3);
\draw(9.25,0.3)node{$\bullet$};\draw[-{Latex[length=3mm]},dashed](9.25,0.3)--(8,.3);
\draw(3.25,0.3)node{$\bullet$};\draw[-{Latex[length=3mm]},dashed](3.25,0.3)--(3.25,-.5)--(10,-.5)--(10,0);

\draw(3.6,2.5)node{bonbon};\draw[-{Latex[length=3mm]},dashed](3.6,2)--(3.6,1);
\end{tikzpicture}

\item Proposition C \mip{bonbon.valeur == bonbon.succ.valeur}


\newpage

\item On complète la fonction \mip{dernier_chaine}.
\begin{minted}[linenos]{python}
def dernier_chaine(n):
    bonbon=creer_collier(n)
    while bonbon.valeur!=bonbon.succ.valeur:
        bonbon.pred.succ=bonbon.succ
        bonbon.succ.pred=bonbon.pred
        bonbon=bonbon.succ.succ.succ
    return bonbon.valeur
\end{minted}

\end{enumerate}




\bigskip

\textbf{Exercice 3} (8 points)

\medskip

\textit{Cet exercice porte sur les bases de données, la programmation en Python, la
récursivité et les algorithmes}

\textit{de parcours de graphes.}

\smallskip

\textbf{Partie A}

\begin{enumerate}
\item La clé \texttt{id\_avion} n'est pas unique dans la table \texttt{reservation} et 
ne peut donc pas servir de clé primaire pour cette table.

\item Le couple \texttt{(id\_vol,id\_passager)} peut servir de clé primaire.

\item Dans une relation, une clé étrangère permet de référencer une autre table (et donc de faire des jointures).

\item La requête \texttt{SELECT id\_vol FROM vol WHERE aeroport\_arr='CDG';} renvoie les \texttt{id\_vol} des vols dont l'aéroport d'arrivée est Paris-Charles de Gaulle. Avec l'extrait de table proposé, on obtient cette table :

\begin{center}
\begin{tabular}{|c|}
\hline
\texttt{id\_vol}\\
\hline
\texttt{AI0015}\\
\texttt{AI0258}\\
\texttt{AI0292}\\
\hline

\end{tabular}
\end{center}

\item La requête \texttt{SELECT aeroport.ville FROM aeroport JOIN vol ON aeroport.id\_aeroport=vol.aeroport\_arr}

\texttt{WHERE vol.aeroport\_dep=CDG ;} convient.

\item On peut mettre à jour la table passager avec la requête

\texttt{UPDATE passager SET d\_total=16 WHERE id\_passager=5; }
\item L'erreur commise ici est une erreur d'intégrité, la clé primaire \texttt{AI0256} est déjà utilisée pour le vol Paris-Sidney, il faut choisir une autre valeur pour la clé primaire, p.ex \texttt{INSERT INTO vol VALUES('AI1024','CDG','YUL',6);}
\end{enumerate}


\textbf{Partie B}

\begin{enumerate}[resume]
\item La valeur de \mip{graphe_airinfo['T']['P']} est \mip{10}.
\item On écrit la fonction \mip{vol_direct} demandée.

\begin{minted}{python}
def vol_direct(graphe,ville1,ville2):
    return ville2 in graphe[ville1]
\end{minted}
\item On écrit la fonction \mip{liste_villes_proches} demandée.
\begin{minted}{python}
def liste_villes_proches(graphe,ville,d_max):
    return [v for v in graphe[ville] if graphe[ville][v]<=d_max]
\end{minted}
\item On construit le graphe représentant le réseau aérien de la compagnie \textit{Droidevant}.

\vspace{-3cm}
\hfill\psset{unit=0.8cm}\begin{pspicture}(6.5,3.5)
\rput(0.5,2){\circlenode{W}{W}}
\rput(3,3){\circlenode{P}{P}}
\rput(2.5,0.5){\circlenode{B}{B}}
\rput(6,.5){\circlenode{T}{T}}
\rput(4.5,3){\circlenode{S}{S}}

\ncarc[arcangle=0]{W}{P}\ncput*{6}
\ncarc[arcangle=0]{W}{B}\ncput*{7}
\ncarc[arcangle=0]{P}{B}\ncput*{1}
\ncarc[arcangle=0]{T}{S}\ncput*{8}
\end{pspicture}

%\vspace{-2cm}
\item Proposition A : le graphe de la compagnie AirInfo est connexe.


\item La fonction \mip{parcours} est récursive car elle s'appelle elle-même.


\item Après l'exécution du code, \mip{visitees1} contient \mip{['W', 'P', 'T', 'B', 'S']} 

et \mip{visitees2} contient \mip{['W', 'P', 'B']}. 

\item Le parcours de graphe utilisé ici est Proposition C : un parcours en profondeur.

\newpage

\item On complète la fonction \mip{est_connexe}.


\begin{minted}[linenos]{python}
def est_connexe(graphe):
    """Vérifie si un graphe est connexe."""
    depart=ville_arbitraire(graphe)
    visitees=[]
    parcours(graphe,visitees,depart)
    return len(visitees)==len(graphe)
\end{minted}

\item L'appel \mip{mystere(graphe_airinfo,'W',[],0,'B')} produit l'affichage ci-dessous.

\begin{minted}{python}
['W', 'P', 'T', 'B'] 25
['W', 'P', 'S', 'T', 'B'] 40
\end{minted}


\item De façon générale, l’appel \mip{mystere(graphe,
ville, [], 0, arrivee)} affiche les chemins sans boucle d'une ville de départ \mip{ville} à une ville d'arrivée \mip{arrivee} dans le réseau représenté par le graphe \mip{graphe}, ainsi que la distance correspondante.
\end{enumerate}



\end{document}